\documentclass[11pt,a4paper]{article}
\usepackage[T1]{fontenc}
\usepackage[utf8]{inputenc}
\usepackage{amsmath,amssymb,amsthm}
\usepackage[margin=1in]{geometry}
\usepackage[colorlinks=true,linkcolor=blue,urlcolor=blue]{hyperref}
\theoremstyle{plain}
\newtheorem{theorem}{Theorem}
\newtheorem{lemma}[theorem]{Lemma}
\newtheorem{proposition}[theorem]{Proposition}
\newtheorem{corollary}[theorem]{Corollary}
\theoremstyle{definition}
\newtheorem{remark}[theorem]{Remark}

\title{The empirical recurrence for OEIS A250427, proved by transfer matrix}
\author{Adrian Perez Fontelles\\ \small Independent researcher}
\date{3 September 2026}
\begin{document}
\maketitle

\begin{abstract}
OEIS A250427 counts arrays over $\{0,1\}$ in which a statistic taken over short runs
of cells produces a derived array that is monotone in a named direction. The entry carries an
empirical recurrence of order $16$ contributed by R.~H.~Hardin. It is true, and it is
decidable rather than empirical: every comparison the condition makes lies inside a bounded
window of consecutive rows, so the count is a walk count on $S=81$ states, and the
conjectured recurrence is then settled by a finite exact computation.
\end{abstract}

\noindent\small 2020 Mathematics Subject Classification. 05A15, 15A18.\normalsize

\section{The conjecture}

OEIS A250427 is ``Number of (n+1)X(3+1) 0..1 arrays with nondecreasing sum of every two consecutive values in every row and column.''. It has offset $1$ and begins
\[
81, 324, 1296, 3600, 10000, 22500, 50625, 99225,\ \dots
\]
The entry states, as a conjecture contributed by its author and never marked settled:
\begin{quote}
Empirical: a(n) = 2*a(n-1) +6*a(n-2) -14*a(n-3) -14*a(n-4) +42*a(n-5) +14*a(n-6) -70*a(n-7) +70*a(n-9) -14*a(n-10) -42*a(n-11) +14*a(n-12) +14*a(n-13) -6*a(n-14) -2*a(n-15) +a(n-16)
\end{quote}
As of the ``Last modified'' line on the live entry (Jul 23 2025 12:23:52, revision 8) this is still
recorded as empirical, and nothing on the entry records it as proved.

\section{What is being counted}

The array has $4$ cells across and entries $x(i,j)\in\{0,1\}$.

Fix the statistic $\sigma(v)=\sum_{t} v_t$ of a run $v$ of $2$ consecutive cells --- in words,
the sum of the $2$ values. Two derived arrays are formed and both must be
nondecreasing in their own direction: reading along a row,
\[
\sigma\bigl(x(i,j),\dots,x(i,j{+}1)\bigr)\ \le\
\sigma\bigl(x(i,j{+}1),\dots,x(i,j{+}2)\bigr),
\]
and reading down a column,
\[
\sigma\bigl(x(i,j),\dots,x(i{+}1,j)\bigr)\ \le\
\sigma\bigl(x(i{+}1,j),\dots,x(i{+}2,j)\bigr).
\]
The same statistic and the same run length serve both directions, so the condition is
unchanged by transposing the array.

Nothing in the condition is invariant under renaming the letters --- each statistic either does
arithmetic on them or compares them by size --- so no reduction to equality patterns is
available, and the walk is run over the alphabet the entry names.

\section{A bounded window}

\begin{lemma}\label{lem:win}
The row comparisons involve one array row each. A column comparison involves the
$2+1=3$ consecutive rows $i,\dots,i{+}2$ and no others.
\end{lemma}

\begin{proof}
The two runs compared start at rows $i$ and $i+1$ and have length $2$, so together they
occupy rows $i$ through $i+2$.
\end{proof}

So the window of the last $2$ rows is a state, each row of it already satisfying its own row
comparisons, and one step appends a row and settles the column comparisons that the new row
completes. With $M$ the adjacency matrix of that digraph and $\iota=\tau=\mathbf{1}$,
\[
a(n)\;=\;\iota^{\!\top}M^{\,j}\tau ,\qquad S=81,
\]
a walk of no steps being an array of exactly $2$ rows.

The walk index $j$ and the entry's own $n$ differ by a constant, read off by matching the
published terms rather than assumed. In particular $a$ is $C$-finite.

\section{The proof}

\begin{lemma}\label{lem:crit}
Let $q(t)=t^{16}-\sum_i c_it^{16-i}$ be the characteristic polynomial of the
conjectured recurrence. The residual $a(n)-\sum_ic_ia(n-i)$ equals
$\iota^{\!\top}M^{\,j}q(M)\tau$ for the corresponding walk index $j$, so the recurrence
holds from some point on if and only if $\iota^{\!\top}M^{j}q(M)\tau=0$ for all large $j$,
and it is enough to examine $j<S$.
\end{lemma}

\begin{proof}
Factor $M^{j}$ out of $M^{j+16}-\sum_ic_iM^{j+16-i}$. For the bound, $M^{S}$
is an integer combination of $I,M,\dots,M^{S-1}$ by Cayley--Hamilton, so the numbers
$u_j=\iota^{\!\top}M^{j}q(M)\tau$ obey the monic recurrence given by the characteristic
polynomial of $M$; $S$ consecutive zeros force every later one, and the last nonzero $u_j$
pins the threshold exactly.
\end{proof}

\begin{theorem}
\[
a(n)\;=\;2\,a(n-1) + 6\,a(n-2) - 14\,a(n-3) - 14\,a(n-4) + 42\,a(n-5) + 14\,a(n-6) - 70\,a(n-7) + 70\,a(n-9) - 14\,a(n-10) - 42\,a(n-11) + 14\,a(n-12) + 14\,a(n-13) - 6\,a(n-14) - 2\,a(n-15) + a(n-16)
\]
for every $n>16$.
\end{theorem}

\begin{proof}
The vector $w=q(M)\tau$ was formed by $16$ matrix--vector products in exact integer
arithmetic and $\iota^{\!\top}M^{j}w$ evaluated for $j=0,1,\dots$, again exactly; those
integers vanish from the index corresponding to $n=16+1$ onwards and the run continued
until the working vector was identically zero, which settles every larger $j$ at once.
\end{proof}

\section{Verification}

Four checks, each independent of the algebra above.

First, the model reproduces all $27$ terms the entry publishes, exactly, in integer
arithmetic.

Second, the count was recomputed for the smallest arrays straight from the definition: fill the
cells in row major order and test each comparison the moment both of its runs are complete,
in the orientation the entry's own name writes. No transfer matrix and no transposition are
involved, and the published terms came back.

Third, the conjectured recurrence was evaluated directly on the published terms, with no
matrices involved, and holds wherever the proved range and the data overlap.

Fourth, the annihilation test of Lemma~\ref{lem:crit} was carried out in exact integer
arithmetic on the whole vector rather than on a sampled prefix.

\begin{thebibliography}{9}
\bibitem{oeis} The OEIS Foundation, \emph{The On-Line Encyclopedia of Integer
Sequences}, \url{https://oeis.org}, 2026. Sequence A250427.
\bibitem{stanley} R.~P.~Stanley, \emph{Enumerative Combinatorics}, Volume 1, 2nd ed.,
Cambridge University Press, 2011. (Transfer-matrix method, Section 4.7.)
\bibitem{horn} R.~A.~Horn and C.~R.~Johnson, \emph{Matrix Analysis}, 2nd ed., Cambridge
University Press, 2013.
\end{thebibliography}

\end{document}
